\subsection{Tangent Spaces}

\lecdate{Lec 11 - Feb 10 (Week 6)}

recall: if $M\subseteq\bR^{n}$ an embedded mfld, the \textbf{tangent space}
at $p\in M$ is $T_{p}M$ given as:
\begin{equation*}
	T_{p}M \ = \ \set{\trm{velocity vectors of curves } \gamma:(a,b)\gto M
	\trm{ at } 0 \trm{ where } \gamma(0)=p} \subseteq \bR^{n}
\end{equation*}
this is in fact an $m$-dim'l vspace.
indeed, $M^{m}$ is locally the img of some $\vphi:U\subseteq\bR^{m}\gto\bR^{n}$,
so $T_{p}M=D\vphi\big\rvert_{\vphi^{-1}(p)}$ where $p\in M$.
note $\vphi$ is not a chart map here.
by chain rule, given $v\in\bR^{n}$, take $\gamma(t)=\vphi^{-1}(p)+tv$. then:
\begin{equation*}
	\frac{\d}{\d t}(\vphi\circ\gamma)\big\rvert_{0}
	\ = \ D\vphi_{p}\left(\frac{\d\gamma}{\d t}\bigg\rvert_{0}\right)
	\ \implies \ D\vphi\big\rvert_{\vphi^{-1}(p)}(\bR^{m})\subseteq T_{p}M
\end{equation*}
to get the reverse inclusion, start with $\gamma:(a,b)\gto M$ and
use $\vphi^{-1}\circ\gamma$. ok idt that was all that important tbh

using directional derivatives, we motivate the defn for an abstract mfld.
\begin{defn}
	a \textbf{tangent vector} at $p\in M$ is a linear map $v:C^{\infty}(M)\gto\bR$
	of the form:
	\begin{equation*}
		v(f) \ = \ \frac{\d}{\d t}\bigg\rvert_{t=0}f(\gamma(t))
	\end{equation*}
	where $\gamma:(a,b)\gto M$ is some smooth curve with $\gamma(0)=p$.
\end{defn}
the idea is to use curves $\gamma$ to take dir derivatives of $f$ at a pt $p$.
in this sense, each fixed $v$ represents a direction along which we can
take the dir derivative of a fn.
this identifies it (implicitly, later explicitly) with a particular vector in $M$
which represents the dir more manifestly.
this distinct identification separates the direction and its behaviour on fns
with the vector in $M$ itself; in this sense, they are really two ways of writing the same thing.
tangent spaces drawn in $\bR^{n}$ are thus ``geometric representations" of them.

remarks:
\begin{enumerate}
	\item this is a subset of the inf-dim $\Hom(C^{\infty}(M),\bR)$
	\item can see it as an equiv rel on curves; that is, $v$ is an equiv class
		of $\gamma$ inducing the same linear map.
		more precisely, curves $\gamma_{1},\gamma_{2}$ are equiv if $\gamma_{1}(0)=\gamma_{2}(0)=p$,
		and $\gamma_{1}'(0)=\gamma_{2}'(0)$. $v$ is then the equiv class, and $v$ intuitively
		represents the vec $\gamma_{i}'(0)$.
\end{enumerate}
in other words, want to think of them not as vectors but as dir derivatives.
we provide an alt coord-wise defn.

\begin{defn}
	let $p\in M,(U,\vphi)$ chart arnd $p$. $v:C^{\infty}(M)\gto\bR$ a
	\textbf{tangent vector} at $p$ if it is a linear combo of functionals in $\Hom(C^{\infty}(M),\bR)$:
	\begin{equation*}
		f \ \mto \ \sum_{i=1}^{m}a_{i}\p_{i}\left(f\circ\vphi^{-1}\big\rvert_{\vphi(p)}\right)
	\end{equation*}
	for some $a\in\bR^{m}$, where $\p_{i}$ is the $i$-th partial deriv.
\end{defn}
here, $a\in\bR^{m}$ represents the dir vector mapped to coords in $\bR^{m}$.
remarks:
\begin{enumerate}
	\item by defn, is a $m$-dim'l linear subspace of $\Hom(C^{\infty}(M),\bR)$
	\item not obv that it doesnt depend on $(U,\vphi)$
\end{enumerate}

\begin{thm}
	the two defn's are equiv.
\end{thm} \

\begin{pf}[source=Primary Source Material]
	given $v$, let $\tilde{\gamma}(t)=\vphi(p)+ta$ a curve $I\sto\vphi(U)$
	and $\gamma=\vphi^{-1}\circ\gamma$. then:
	\begin{equation*}
		v(f)
		\ = \ \sum_{i=1}^{m}a_{i}\p_{i}\left(f\circ\vphi^{-1}\big\rvert_{\vphi(p)}\right)
		\ = \ \frac{\d}{\d t}\bigg\rvert_{t=0}(f\circ\vphi^{-1})(\tilde{\gamma}(t))
		\ = \ \frac{\d}{\d t}\bigg\rvert_{t=0}f(\gamma(t))
	\end{equation*}
	by chain rule, so $v$ of the [first] form.
	now, if $v$ [of the first form] for some $\gamma$ w $\gamma(0)=p$, then:
	\begin{equation*}
		v(f) \ = \ \frac{\d}{\d t}\bigg\rvert_{t=0}(f\circ\vphi^{-1})(\vphi\circ\gamma(t))
		\ = \ \nabla(f\circ\vphi^{-1})\bigg\rvert_{\vphi(p)}\cdot\frac{\d}{\d t}(\vphi\circ\gamma)\bigg\rvert_{t=0}
	\end{equation*}
	letting $a:=\dfrac{\d}{\d t}(\vphi\circ\gamma)\bigg\rvert_{t=0}$,
	this is then of the second form.
\end{pf}

ah. so, here we're identifying $v=\sum a_{i}\p_{i}$, or in other words,
vectors $a$ with ``operators" $v$. thus, $v$ represents a directional deriv
in the direction of $a$. in the first defn, $a$ is given as the derivative
of $\gamma$ at $t=0$. $<$- previous ramblings

ok now onto the txtbk. this is gonna suck.

using the 2nd defn, any chart $U$ arnd $p$ clearly gives an iso $T_{p}M\simeq\bR^{m}$
by identifying $v$ with $a$.
specifically, if $U\subseteq M$ open, then $T_{p}U\subseteq C^{\infty}(U)$ is
the subspace spanned by the partials at $p$:
\begin{equation*}
	\frac{\p}{\p x_{1}}\bigg\rvert_{p}, \ \dots \ , \frac{\p}{\p x_{m}}\bigg\rvert_{p}
\end{equation*}
this basis identifies $T_{p}U\simeq\bR^{m}$.
coords $a_{i}$ can then be obtained by taking $v(\pi_{i})$.

there is yet another characterization which is coord-free and doesn't rely on curves.
first, note that clearly every $v$ satisfies the product rule, also known as
the Leibniz rule:
\begin{equation*}
	v(fg) \ = \ v(f)g(p) + f(p)v(g)
\end{equation*}
evidently, there are many $v\in\Hom(C^{\infty}(M),\bR)$ which do not satisfy
this property. however, we claim the following:
\begin{thm}
	a linear $v:C^{\infty}(M)\gto\bR$ defines an elem of $T_{p}M$ iff it
	satisfies the product rule.
\end{thm}
the pf is rly long and uses ``Hadamard's Lemma", i highly doubt its important.
im skipping it. anyway, it gives us this third defn:

\begin{defn}
	given $v:C^{\infty}(M)\gto\bR$, $v\in T_{p}M$ iff:
	\begin{equation*}
		v(fg) \ = \ v(f)g(p) + f(p)v(g)
	\end{equation*}
	for all $f,g$.
\end{defn}
there is \textit{yet another}, which seemingly involves dual spaces.
the book warns to skip it tho, esp if you dont like abstractions.
ill read it, but im not taking notes on it.

notice that ``velocity vectors" are naturally elements of the tangent space.
given a curve $\gamma$, then $\gamma'(t_{0})\in T_{\gamma(t_{0})}M$ is given by:
\begin{equation*}
	(\gamma'(t_{0}))(f) \ = \ \frac{\d}{\d t}\bigg\rvert_{t_{0}}f(\gamma(t))
\end{equation*}
